% ============================================================
\chapter{Interest Rate Models}
\label{ch:interest_rates}
% ============================================================

Until now, we have treated the interest rate as a constant. But anyone who has observed bond markets knows that rates are anything but constant: they form a \emph{curve} --- the term structure --- that evolves stochastically over time. Modelling this evolution is one of the richest challenges in quantitative finance. Oldrich Vasicek, in 1977, proposed the first stochastic short-rate model with mean reversion, drawing inspiration from the Ornstein--Uhlenbeck process. John Cox, Jonathan Ingersoll, and Stephen Ross (1985) corrected Vasicek's flaw --- the possibility of negative rates --- by introducing a volatility proportional to $\sqrt{r}$. Then David Heath, Robert Jarrow, and Andrew Morton (1992) reversed the perspective entirely: instead of modelling the short rate, they modelled the entire forward curve directly. Each approach serves different needs, and understanding them is understanding how markets price time and uncertainty.

\section{Term Structure of Interest Rates}

\begin{definition}[Zero-Coupon Bond and Yield Curve]
A \textbf{zero-coupon bond} with maturity $T$ is a contract paying 1 unit
at time $T$. Its price at time $t$ is denoted $P(t,T)$. The
\textbf{continuously compounded zero rate} is:
\[
    R(t,T) = -\frac{\ln P(t,T)}{T - t}.
\]
The curve $T \mapsto R(t,T)$ is the \textbf{term structure} of interest
rates at time $t$.
\end{definition}

\begin{definition}[Instantaneous Forward Rate]
The \textbf{instantaneous forward rate} is:
\[
    f(t,T) = -\frac{\partial \ln P(t,T)}{\partial T},
\]
so that $P(t,T) = \exp\!\left(-\int_t^T f(t,u)\,du\right)$.
\end{definition}

\begin{definition}[Short Rate]
The \textbf{short rate} (or instantaneous rate) is:
$r_t = f(t,t) = \lim_{T \to t^+} R(t,T)$.
\end{definition}

\section{Short-Rate Models}

\subsection{Vasicek Model (1977)}

\begin{definition}[Vasicek Model]
The short rate follows an Ornstein--Uhlenbeck process:
\[
    dr_t = a(b - r_t)\,dt + \sigma\,dW_t,
\]
with $a > 0$ (mean-reversion speed), $b$ (long-term rate),
$\sigma > 0$ (volatility).
\end{definition}

\begin{theorem}[Zero-Coupon Price --- Vasicek]
\label{thm:vasicek_zc}
The zero-coupon bond price has the affine form:
\[
    P(t,T) = A(t,T)\exp\!\bigl(-B(t,T)\,r_t\bigr),
\]
where, with $\tau = T - t$:
\begin{align}
    B(\tau) &= \frac{1 - e^{-a\tau}}{a}, \\
    A(\tau) &= \exp\!\left[\left(b - \frac{\sigma^2}{2a^2}\right)
    \bigl(B(\tau) - \tau\bigr) - \frac{\sigma^2}{4a}B(\tau)^2\right].
\end{align}
\end{theorem}

\begin{remark}
The Vasicek model allows negative interest rates, which was considered a
defect until the 2010s but proved realistic in the European negative rate
environment.
\end{remark}

\subsection{Cox--Ingersoll--Ross Model (CIR, 1985)}

\begin{definition}[CIR Model]
\[
    dr_t = a(b - r_t)\,dt + \sigma\sqrt{r_t}\,dW_t.
\]
The $\sqrt{r_t}$ term in the diffusion ensures $r_t \geq 0$ under the
Feller condition $2ab \geq \sigma^2$.
\end{definition}

\begin{theorem}[Zero-Coupon Price --- CIR]
The price is affine in $r_t$:
$P(t,T) = A(\tau)\exp(-B(\tau)r_t)$ with:
\begin{align}
    B(\tau) &= \frac{2(e^{\gamma\tau} - 1)}
    {(\gamma + a)(e^{\gamma\tau} - 1) + 2\gamma}, \\
    A(\tau) &= \left[\frac{2\gamma e^{(a+\gamma)\tau/2}}
    {(\gamma+a)(e^{\gamma\tau}-1) + 2\gamma}\right]^{2ab/\sigma^2},
\end{align}
where $\gamma = \sqrt{a^2 + 2\sigma^2}$.
\end{theorem}

\subsection{Hull--White Model (1990)}

\begin{definition}[Hull--White Model]
An extension of Vasicek with time-dependent parameters:
\[
    dr_t = \bigl[\theta(t) - a\,r_t\bigr]\,dt + \sigma\,dW_t,
\]
where $\theta(t)$ is chosen to calibrate the initial yield curve exactly:
\[
    \theta(t) = \frac{\partial f^M(0,t)}{\partial t} + a\,f^M(0,t)
    + \frac{\sigma^2}{2a}(1 - e^{-2at}),
\]
with $f^M(0,t)$ the market forward rate.
\end{definition}

\begin{codeblock}[title=\textbf{Vasicek and CIR Models}]
\begin{minted}{python}
import numpy as np

def vasicek_bond(r, t, T, a, b, sigma):
    """Vasicek zero-coupon price."""
    tau = T - t
    B = (1 - np.exp(-a * tau)) / a
    A = np.exp((b - sigma**2/(2*a**2)) * (B - tau)
               - sigma**2/(4*a) * B**2)
    return A * np.exp(-B * r)

def cir_bond(r, t, T, a, b, sigma):
    """CIR zero-coupon price."""
    tau = T - t
    gamma = np.sqrt(a**2 + 2*sigma**2)
    eg = np.exp(gamma * tau)
    denom = (gamma + a)*(eg - 1) + 2*gamma
    B = 2*(eg - 1) / denom
    A = (2*gamma * np.exp((a + gamma)*tau/2) / denom)**(2*a*b/sigma**2)
    return A * np.exp(-B * r)

# Yield curve
r0 = 0.03
maturities = np.linspace(0.25, 30, 100)
P_vas = [vasicek_bond(r0, 0, T, 0.5, 0.04, 0.01) for T in maturities]
P_cir = [cir_bond(r0, 0, T, 0.5, 0.04, 0.06) for T in maturities]
R_vas = -np.log(P_vas) / maturities
R_cir = -np.log(P_cir) / maturities

print(f"Vasicek R(10y): {R_vas[39]:.4f}")
print(f"CIR     R(10y): {R_cir[39]:.4f}")
\end{minted}
\end{codeblock}

\section{Heath--Jarrow--Morton (HJM) Framework}

\begin{theorem}[HJM Framework (1992)]
The HJM framework models instantaneous forward rates directly:
\[
    df(t,T) = \alpha(t,T)\,dt + \sigma(t,T)\,dW_t.
\]
The no-arbitrage condition imposes the \textbf{HJM drift condition}:
\[
    \alpha(t,T) = \sigma(t,T)\int_t^T \sigma(t,u)\,du.
\]
\end{theorem}

\begin{remark}
The HJM framework is very general but typically leads to non-Markovian
models (except for specific choices of $\sigma$). The Vasicek and Hull--White
models are special cases of HJM.
\end{remark}

\section{Interest Rate Product Pricing}

\subsection{Caps and Floors}

\begin{definition}[Caplet]
A \textbf{caplet} on the period $[T_i, T_{i+1}]$ with strike rate $K$
and notional $N$ pays:
\[
    N\,\delta_i\,(L(T_i, T_i, T_{i+1}) - K)^+,
\]
where $L(T_i, T_i, T_{i+1})$ is the LIBOR rate and $\delta_i = T_{i+1} - T_i$.
\end{definition}

\begin{theorem}[Black's Formula for Caplets]
Under the Black model, the caplet price is:
\[
    \text{Caplet} = N\,\delta_i\,P(0,T_{i+1})
    \bigl[L_0\Phi(d_1) - K\Phi(d_2)\bigr],
\]
with $d_1 = \frac{\ln(L_0/K) + \sigma_i^2 T_i/2}{\sigma_i\sqrt{T_i}}$,
$L_0 = L(0, T_i, T_{i+1})$.
\end{theorem}

\subsection{Swaptions}

\begin{definition}[Swaption]
A \textbf{swaption} gives the right to enter into an interest rate swap at
a future date. A payer swaption with strike $K$ and maturity $T_0$ on a
swap with dates $T_0, \ldots, T_n$ has payoff:
\[
    \left(\sum_{i=0}^{n-1}\delta_i P(T_0, T_{i+1})\right)
    (S_{T_0} - K)^+,
\]
where $S_{T_0}$ is the par swap rate.
\end{definition}

\begin{codeblock}[title=\textbf{Short Rate Path Simulation (Vasicek)}]
\begin{minted}{python}
def vasicek_paths(r0, a, b, sigma, T, N, M):
    """Simulate M short rate paths under Vasicek."""
    dt = T / N
    r = np.zeros((M, N + 1))
    r[:, 0] = r0
    for i in range(N):
        dW = np.sqrt(dt) * np.random.standard_normal(M)
        r[:, i+1] = r[:, i] + a*(b - r[:, i])*dt + sigma*dW
    return r

def cir_paths(r0, a, b, sigma, T, N, M):
    """Simulate M short rate paths under CIR (truncated Euler)."""
    dt = T / N
    r = np.zeros((M, N + 1))
    r[:, 0] = r0
    for i in range(N):
        dW = np.sqrt(dt) * np.random.standard_normal(M)
        r_pos = np.maximum(r[:, i], 0)
        r[:, i+1] = (r[:, i] + a*(b - r_pos)*dt
                     + sigma*np.sqrt(r_pos)*dW)
        r[:, i+1] = np.maximum(r[:, i+1], 0)
    return r

np.random.seed(42)
r_vas = vasicek_paths(0.03, 0.5, 0.04, 0.01, 10, 2520, 5)
r_cir = cir_paths(0.03, 0.5, 0.04, 0.06, 10, 2520, 5)
\end{minted}
\end{codeblock}

\section{LIBOR Market Model (BGM/LMM)}

\begin{definition}[BGM/LMM Model]
The \textbf{LIBOR Market Model} (Brace--Gatarek--Musiela) models forward
LIBOR rates $L_i(t) = L(t, T_i, T_{i+1})$ directly under the $T_{i+1}$
forward measure:
\[
    dL_i(t) = \sigma_i(t)L_i(t)\,dW_t^{T_{i+1}}.
\]
Under other measures, a drift term appears depending on correlations
between forward rates.
\end{definition}

\section{Exercises}

\begin{exercise}[Yield Curve]
\begin{enumerate}
    \item Compute and plot the yield curve $T \mapsto R(0,T)$ for Vasicek
          and CIR with $r_0 = 0.03$, $a = 0.5$, $b = 0.04$,
          $\sigma_V = 0.01$, $\sigma_C = 0.06$.
    \item Compare curve shapes (upward-sloping, downward-sloping, humped).
    \item Study the impact of parameters $a$, $b$, $\sigma$ on the shape.
\end{enumerate}
\end{exercise}

\begin{exercise}[Monte Carlo Bond Pricing]
\begin{enumerate}
    \item Price a 5-year zero-coupon bond by Monte Carlo under Vasicek:
          $P(0,T) = \E^\mathbb{Q}\!\left[e^{-\int_0^T r_t\,dt}\right]$.
    \item Compare with the analytical formula.
    \item Repeat for the CIR model.
\end{enumerate}
\end{exercise}

\begin{exercise}[Caplet Pricing]
\begin{enumerate}
    \item Price a 1-year caplet using Black's formula.
    \item Price the same caplet by Monte Carlo under the Hull--White model.
    \item Compare the results.
\end{enumerate}
\end{exercise}

\begin{exercise}[Hull--White Calibration]
\begin{enumerate}
    \item Calibrate the Hull--White model to a given yield curve by
          determining $\theta(t)$.
    \item Verify that the model zero-coupon prices match market prices.
\end{enumerate}
\end{exercise}
